% TEXTO DEL RESUMEN (Aseguramos justificación y sin división de palabras)
% Si usaste \hyphenpenalty=10000 globalmente, ya aplica aquí. 
% Si no, puedes añadirlo aquí dentro de un grupo si es necesario.
\justifying % Mantiene el bloque de texto justificado.

El presente Trabajo Especial de Grado tuvo como objetivo principal diseñar un modelo de gestión integral para la optimización de los procesos logísticos en la industria petrolera del estado Anzoátegui. La investigación se enmarcó en un diseño de campo de carácter descriptivo, con una población de 50 especialistas y una muestra de 20. Los resultados revelaron deficiencias significativas en la sincronización de inventarios y transporte, lo que genera retrasos del 15\% en la cadena de suministro. Se concluye que la implementación del modelo propuesto permitirá reducir los tiempos de espera y mejorar la eficiencia operativa en al menos un 10\%, aportando una solución viable a la problemática identificada.
\vspace{1em} % Espacio vertical de una línea, aproximadamente

% 3. PALABRAS CLAVE
\noindent\textbf{Palabras clave:} Logística, Gestión Integral, Optimización, Cadena de Suministro, UNEFA.
